\documentclass[aspectratio=169,10pt,spanish]{beamer}

\input{../../../_preambulo_beamer.tex}
\input{../../../_comandos_beamer.tex}

\selectlanguage{spanish}

\title{MA1001B - Modelación Estadística para la Toma de Decisiones}
\subtitle{Lección 22: Aplicaciones normales: tolerancias e intervalos operativos}
\author{}
\institute{Tecnol\'ogico de Monterrey}
\date{}

\begin{document}

\begin{frame}[plain]
  \vspace{1.2cm}
  {\Large\bfseries MA1001B - Modelación Estadística para la Toma de Decisiones\par}
  \vspace{0.7cm}
  {\large Lección 22: Aplicaciones normales y tolerancias\par}
  \vspace{0.7cm}
  {\color{TecRojo}\rule{\textwidth}{1.2pt}}
  \vspace{0.8cm}

  Facultad MA1001B\\[0.3cm]
  Periodo\\[0.3cm]
  Tecnol\'ogico de Monterrey
\end{frame}

\begin{frame}{La pregunta de especificación}
  \begin{alertblock}{Pregunta guía}
    Si el peso de llenado es $N(50,4)$ gramos y la especificación es
    $[44,56]$, ¿qué fracción de unidades está en especificación, y qué ocurre
    si se mueve la media?
  \end{alertblock}

  \vspace{0.6em}
  Al final de esta lección, usted puede:
  \begin{itemize}
    \item convertir límites de especificación en cortes $z$;
    \item calcular el rendimiento como $P(\mathrm{LIE}<X<\mathrm{LSE})$;
    \item cuantificar la caída de rendimiento tras un desplazamiento de la media;
    \item explicar por qué un rendimiento alto no es lo mismo que estar centrado.
  \end{itemize}

  \vfill
  \scriptsize Todos los pesos de llenado usados hoy son sintéticos. No se usan datos reales de empresa.
\end{frame}

\begin{frame}{Proceso centrado y cortes de especificación}
  \small
  Modelo $X\sim N(50,4^2)$. Especificación $[44,56]$:
  \[
    z_{\mathrm{LIE}}=\frac{44-50}{4}=-1.5,\qquad
    z_{\mathrm{LSE}}=\frac{56-50}{4}=1.5.
  \]
  Rendimiento en especificación:
  \[
    P(44<X<56)=0.866386.
  \]
  Cada cola es $0.066807$ porque el proceso está centrado.

  \vspace{0.3em}
  \begin{exampleblock}{Línea de llenado sintética}
    La especificación es un intervalo del cliente. El rendimiento es el área
    entre los cortes.
  \end{exampleblock}
\end{frame}

\begin{frame}[fragile]{Paso 1: Rendimiento en especificación}
  \begin{columns}[T]
    \begin{column}{0.42\textwidth}
      \small
      \begin{lstlisting}
z_low = (44 - 50) / 4
z_high = (56 - 50) / 4
yield_p = (
    model.cdf(56)
    - model.cdf(44)
)
      \end{lstlisting}

      \begin{block}{Salida verificada}
        $z=\pm 1.500000$\\
        Cada cola $=0.066807$\\
        Rendimiento $=0.866386$
      \end{block}
    \end{column}
    \begin{column}{0.58\textwidth}
      \centering
      \includegraphics[width=\textwidth]{../figuras/tolerancias_normal_01_rendimiento_espec.png}
      \vspace{0.2em}
      \scriptsize Rendimiento centrado del Paso 1
    \end{column}
  \end{columns}
\end{frame}

\begin{frame}{Un desplazamiento de la media mueve los cortes $z$}
  \small
  La especificación permanece en $[44,56]$. Después de $\mu=52$,
  \[
    z_{\mathrm{LIE}}=\frac{44-52}{4}=-2,\qquad
    z_{\mathrm{LSE}}=\frac{56-52}{4}=1.
  \]
  El rendimiento cae a $0.818595$. La caída frente al proceso centrado es
  $0.047791$. La cola superior crece a $0.158655$.
\end{frame}

\begin{frame}[fragile]{Paso 2: Rendimiento después de $\mu=52$}
  \begin{columns}[T]
    \begin{column}{0.40\textwidth}
      \small
      \begin{lstlisting}
shifted = norm(52, 4)
new_yield = (
    shifted.cdf(56)
    - shifted.cdf(44)
)
      \end{lstlisting}

      \begin{block}{Salida verificada}
        Cortes $z$ $-2$ y $1$\\
        Rendimiento $=0.818595$\\
        Caída $=0.047791$
      \end{block}
    \end{column}
    \begin{column}{0.60\textwidth}
      \centering
      \includegraphics[width=\textwidth]{../figuras/tolerancias_normal_02_desplazamiento_media.png}
      \vspace{0.2em}
      \scriptsize Densidades centrada y desplazada del Paso 2
    \end{column}
  \end{columns}
\end{frame}

\begin{frame}{Paso 3: Un rendimiento alto no implica estar centrado}
  \begin{columns}[T]
    \begin{column}{0.42\textwidth}
      \small
      Reloj estrecho descentrado: $\mu=52$, $\sigma=2$.

      \vspace{0.4em}
      Cortes $z$ $-4$ y $2$.\\
      Rendimiento $=0.977218$, más alto que el centrado $0.866386$.

      \vspace{0.4em}
      \textbf{Límite:} la conformidad con la especificación no es lo mismo que
      estar centrado. La Lección 23 estudia la distribución muestral de
      $\bar x$.
    \end{column}
    \begin{column}{0.58\textwidth}
      \centering
      \includegraphics[width=\textwidth]{../figuras/tolerancias_normal_03_centrado_vs_conforme.png}
      \vspace{0.2em}
      \scriptsize Centrado-ancho frente a estrecho descentrado del Paso 3
    \end{column}
  \end{columns}
\end{frame}

\begin{frame}{Verificación de conceptos}
  \begin{enumerate}
    \item Calcule los cortes $z$ para la especificación $[44,56]$ en $N(50,4)$.
    \item ¿Por qué las dos colas igualan $0.066807$ en el Paso 1?
    \item Después de $\mu=52$, ¿cuál cola crece, y a qué valor?
    \item ¿Por qué el rendimiento $0.977218$ puede seguir siendo la historia operativa equivocada?
  \end{enumerate}

  \vspace{0.6em}
  \begin{block}{Conexión con la Lección 23}
    La sesión puede continuar directamente con la distribución muestral de la
    media muestral y su error estándar.
  \end{block}

  \vfill
  \scriptsize Fuente: OpenStax, \textit{Introductory Business Statistics 2e},
  Capítulo 6, Sección 6.2. CC BY-NC-SA.
\end{frame}

\appendix



\end{document}
